% section: related
\section{Related work}
\label{sec:related}

\paragraph{Chart question answering is close to saturated; chart grounding is not.}
This asymmetry is the premise of the present work. A current 2B model scores \result{79.1} on
ChartQA with no task-specific training at all, and the best published improvement over a comparable
prompted 3B model is $+3.8$ points, obtained on four H100s. Chart \emph{grounding} sits far lower:
the strongest published RefChartQA human-subset AP@0.5 is \result{32.83}, produced by a 3B model
trained for a single epoch with boxes emitted as plain text. Independent evidence points the same
way, with a current model scoring roughly 8--15\% on multi-target chart grounding.

We therefore report both targets, and design so that the defensible measured improvement does not
depend on the saturated one.

\paragraph{RefChartQA} supplies the grounding benchmark, its evaluator, and the reference figure this
work targets. It also supplies a directly relevant measurement: adding box output cost a 3B model
$88.80 \rightarrow 84.80$ in answer accuracy, one of two published estimates of the price of
structured output.

\paragraph{Chart-RVR} supplies the other estimate --- chain-of-thought cost $82.0 \rightarrow 73.1$
on a 3B model --- and establishes both that ChartQA improvements are small and expensive, and that
structured output is not free.

\paragraph{START} demonstrates that chart-element spatial supervision improves chart understanding
across four benchmarks. It validates the training signal used here while leaving the grounding
target uncontested, since it does not report RefChartQA AP@0.5.

\paragraph{CURV} contributes a deterministic chart generator and, more usefully for this work, a
negative result: explicit multi-turn grounding at inference time can \emph{lower} accuracy through
cumulative box error and reasoning--grounding mismatch. That finding is the reason the crop stage in
this system is demoted to a reported ablation rather than placed on the critical path.

\paragraph{ChartPoint} supports both the use of box supervision and the coordinate format, finding
that integer $0$--$999$ coordinates outperform decimal ones. Its recipe cost roughly 262 A100-hours
and is not adopted.

\paragraph{Chart-FR1} reports 91.0 on ChartQA at 7B using eight H100s and a GPT-5 teacher. Its
repository contains only a README and a licence, so it is context and explicitly not a reproducible
baseline.

\paragraph{ChartQAPro} shows how brittle chart models are outside their training distribution, with
a strong commercial model falling from 90.5 on ChartQA to 55.81. It is used here, if at all, as a
sealed stress test.
